\chapter{Introduction and Goals}

\section{Introduction}

This document intends to describe the low level design of the SUM_LNM module. The scope of this document is to capture the implementation based on the requirements as defined in 
SDV-8002 AUTOSAR Standard Utility Module Specification. For a better process and higher quality design, this can be seen as an extension to the before mentioned document.

\subsection{System Purpose}

This document defines a set of standard software compositions which are referred to as GM Standardized Utility Modules (SUM). 
Use of the SUM m with respect to the interface to the software infrastructure (BSW Services). This allows functional software to be completely abstracted from 
the software infrastructure. SUM will be used to communicate a set of standardized implementations of the interactions between SW-C and BSW.
The Standard Utility Module, LIN Network Manager (SUM\_LNM), is a software composition located in the Application Layer of the standardized architecture of AUTOSAR.

Main Features:
\begin{itemize}
	\item Tracks and controls the activation status (Full Communication or No Communication) of the LIN Network using SUM\_LNM interfaces and BSW services
	\item Track and switch the LIN Schedule using SUM_LNM interfaces and BSW services
	\item Collect all conditional logic related to switching the LIN Schedule(s) for each LIN Network in which the ECU is the LIN Commander
	\item Execute the logic to switch the LIN Schedules for each LIN Network in which the ECU is the LIN Commander
	\item Enable the LIN Network when enabled by SUM\_SUSD
	\item Perform the sequence necessary to put the LIN Network to sleep and notify SUM\_SUSD
	\item Detect U codes for Node Loss detection of LIN Responders to SUM\_ERRH
	\item Support Diagnostics of Functional Capability via notification to SUM\_ERRH
	\item Interface with Dcm to receive any required DID requests relevant to the MSCL
	\item Interface with Dcm to receive any required RID requests relevant to switching the LIN schedule
	\item Calibration Download Support Feature
	\item Responder Requested Wakeup Handling
	\item Maintain MSCL Information (Status of nodes on the bus)
\end{itemize}

\subsection{Document Purpose}

This document introduces the key ideas of the architecture of SUM_LNM. Its intended audience comprises System Architects, Product Architects, Developers and Testers 
who need a rough understanding on the concepts of the product. 

The document describes the product and its functionality on the component level, the interaction and communication between those components and the integration of 
ETAS wide shared components.

The document does not explain the components design in detail.

\section{Goals of the Architecture}

\subsection{Goals}

The main goals of the architecture are:

\begin{enumerate}
\item \textbf{Configurability}: The module has to be designed with the highest possible level of configurability, in order to be easy to adapt it in the future to changes on the standard or even in some project requirements.
\item \textbf{Maintainability}: The architecture of the component shall be clearly defined and easy to understand, in order to guarantee the maximum possible level of maintainability.
\item \textbf{Performance}: The component shall be designed in order to maximize run-time execution speed as well as minimize memory consumption.
\item \textbf{Correctness of Implementation}: The Detailed Design must to be aligned to the Architecture document (SDV-8002).
\end{enumerate}

\section{Stakeholders}

Stakeholders for SUM_ERRH module are GM, Bosch, Systems engineers, Software architect, QA personnel.

\section{Definitions, Acronyms and Abbreviations}

\input{generated/tables/definitions.tex}

\section{Referenced Documents}

\input{generated/tables/references.tex}